% !TeX root = CV_Tramarin.tex

%%%%%%%%%% ITALIANO

\langit{

\cvitem{}{L'attuale carico didattico istituzionale è di 21 CFU (182 ore), suddivisi in 5 insegnamenti su tre sedi universitarie (Università degli Studi di Modena e Reggio Emilia, Università degli Studi di Bologna, Università degli Studi di Verona).\medskip
    \begin{itemize}
        \item 6 CFU (48 ore) di ``Metodi, tecniche di misura e sensori'' (Laurea interateneo in Ingegneria dei sistemi medicali per la persona, UniVR)
        \item 1 CFU (16 ore) di ``Laboratorio di Elettronica Circuitale'' (Laurea Triennale in Ingegneria Elettronica, Unimore)
        \item 9 CFU (72 ore) di ``Industrial Measurements'' (Laurea Magistrale in Electronics Engineering, Unimore)
        \item 3 CFU (30 ore) di ``Advanced Sensors for Electric Vehicles'' (Laurea Magistrale in Electric Vehicle Engineering, UniBO)
        \item 2 CFU (16 ore) di ``Hardware and Software Co-Design'' (Laurea Magistrale in Electronics Engineering, Unimore)
    \end{itemize}
    \mbox{}\newline
    A questi si aggiunge la titolarità di un insegnamento, per un totale di 10 ore di didattica frontale, attivo dal corrente anno accademico, all'inerno dei ``Percorsi formativi abilitanti per la scuola secondaria'' svolti presso l'ateneo di Modena.\medskip\newline
    Di seguito si trovano i dettagli degli insegnamenti  (o moduli) di cui si ha responsabilità, e le informazioni relative agli insegnamenti con titolarità assunti prima del passaggio al ruolo di Professore di II fascia.}\medskip


\cventry{2024 - presente}{``Hardware and Software Co-Design''}
{Laurea Magistrale in Electronics Engineering, I anno - 1° semestre}
{\newline  Dip. di Ingegneria “Enzo Ferrari'', Università degli Studi di Modena\newline \textbf{Responsabilità di Modulo 2 CFU (16 ore, IMIS-01/B - ex ING-INF/07)} - Titolarità del corso: prof. L. Rovati (SSD: ING-INF/07 e ING-INF/01)\newline
    L'insegnamento, caratterizzante per il primo anno della magistrale in Electronics Engineering, per un totale di 12 CFU, è suddiviso in due moduli di 6 CFU, il primo di area elettronica (ING-INF/01) e il secondo nell'area delle misure elettroniche (ING-INF/07). Questo secondo modulo di 6 CFU è suddiviso in 3 sottomoduli di 2 CFU ciascuno. Nel modulo di 2 CFU di cui si ha titolarità gli studenti vengono introdotti alla simulazione di semplici circuiti elettronici, tramite SPICE, con particolare riferimenti agli stadi di filtraggio del segnale utili in fase di condizionamento del segnale}{\newline 1 anno, dal A.A. 2024/2025}{\href{https://unimore.coursecatalogue.cineca.it/insegnamenti/2024/28528/2024/10001/10858?coorte=2024&schemaid=21692}{Collegamento al catalogo universitario degli insegnamenti}}

\cventry{2025 - }{``Generazione, acquisizione, elaborazione, e distribuzione dei segnali''}{Percorsi formativi abilitanti per la scuola secondaria}
{\newline  Dip. di Ingegneria “Enzo Ferrari'', Università degli Studi di Modena\newline \textbf{Titolarità dell'Insegnamento - 10 ore (SSD: ING-INF/07)}\newline
    L'insegnamento, di 10 ore di didattica frontale, è parte del percorso formativo abilitante per la scuola secondaria, si svolge in presenza, ed è diretto alla classe di insegnamento A040 (Scienze e Tecnologie Elettriche ed Elettroniche), e B015 (Laboratori di Scienze e Tecnologie Elettriche ed Elettroniche)}{\newline 1 anno, A.A. 2024/2025}{\href{https://www.unimore.it/it/didattica/formazione-insegnanti/percorsi-formativi-abilitanti-60-cfu-la-scuola-secondaria}{https://www.unimore.it/it/didattica/formazione-insegnanti/percorsi-formativi-abilitanti-60-cfu-la-scuola-secondaria}}

\cventry{2023 - 2024}{``Collegamenti e reti, RS-232, LAN, Fondamenti sui bus di campo''}{Percorsi formativi abilitanti per la scuola secondaria}
{\newline  Dip. di Ingegneria “Enzo Ferrari'', Università degli Studi di Modena\newline \textbf{Titolarità dell'Insegnamento - 10 ore (SSD: ING-INF/07)}\newline
    L'insegnamento, di 10 ore di didattica frontale, è parte del percorso formativo abilitante per la scuola secondaria, si svolge in modalità blended, ed è diretto alla classe di insegnamento B015 (Laboratori di Scienze e Tecnologie Elettriche ed Elettroniche)}{\newline 1 anno, A.A. 2023/2024}{\href{https://www.unimore.it/it/didattica/formazione-insegnanti/percorsi-formativi-abilitanti-60-cfu-la-scuola-secondaria}{https://www.unimore.it/it/didattica/formazione-insegnanti/percorsi-formativi-abilitanti-60-cfu-la-scuola-secondaria}}

\cventry{2023 - presente}{``Metodi, tecniche di misura e sensori''}
{Laurea interateneo in Ingegneria dei sistemi medicali per la persona, III anno - 1° semestre}
{\newline  Dipartimento di Ingegneria per la medicina di innovazione, Università degli Studi di Verona}{\newline \textbf{Titolarità dell'Insegnamento - 6 CFU (48 ore, IMIS-01/B - ex ING-INF/07)} -\newline
    L'insegnamento si svolge all'interno di un corso di laurea interateneo tra gli atenei di Verona, Modena e Reggio Emilia e Trento. La sede amministrativa del CdS è presso l'Università di Verona, dove vengono svolte le lezioni frontali.\newline 2 anni, dal A.A. 2023/2024}{\href{https://www.corsi.univr.it/?ent=cs&aa=2024\%2F2025&codiceCs=S25&codins=4S009924&discr=&discrCd=&id=1001&menu=studiare&tab=insegnamenti}{Collegamento al catalogo universitario degli insegnamenti}}

\cventry{2023 - presente}{``Laboratorio di Elettronica Circuitale''}
{Laurea Triennale in Ingegneria Elettronica, I anno - 2° semestre}
{\newline  Dip. di Ingegneria “Enzo Ferrari'', Università degli Studi di Modena}{\newline \textbf{Responsabilità di Modulo 1 CFU (16 ore, IMIS-01/B - ex ING-INF/07)} - Titolarità del corso: prof. L. Vincetti (ex ING-INF/01)\newline
    L'insegnamento, orientato fortemente alle attività di laboratorio, è strutturato in 6 CFU suddivisi in moduli di 1 CFU, corrispondenti a 16 ore di attività di laboratorio, dovute alla necessità di suddividere gli studenti in turni consecutivi. Ogni modulo è tenuto da docenti diversi, afferenti ai settori dell'elettronica e delle misure elettroniche. Lo scopo è di fornire un allineamento alle conoscenze di base necessarie per affrontare il corso di Laurea in Ingegneria Elettronica. Il modulo di 1 CFU di cui si ha responsabilità è dedicato all'analisi e progettazione di circuiti elettronici, con particolare attenzione alla progettazione e realizzazione di un sistema di misura. \newline 2 anni, dal A.A. 2023/2024}{\href{https://unimore.coursecatalogue.cineca.it/insegnamenti/2024/27909/2021/10000/10359?coorte=2024}{Collegamento al catalogo universitario degli insegnamenti}}

\cventry{2020 - presente}{``Industrial Measurements''}
{Laurea Magistrale in Electronics Engineering, I anno - 2° semestre}
{\newline  Dip. di Ingegneria “Enzo Ferrari'', Università degli Studi di Modena}{\newline \textbf{Titolarità dell'Insegnamento - 9 CFU (72 ore, IMIS-01/B - ex ING-INF/07)}\newline
    L'insegnamento è strutturato in 6 CFU di lezioni frontali e 3 CFU di laboratorio. Il laboratorio prevede una parte di allineamento all'utilizzo di LabVIEW e una parte di sviluppo autonomo di un'attività progettuale per la realizzazione sistema di misura, basato sulle tecniche e i sensori descritti a lezione.\newline 5 anni, dal A.A. 2020/2021}{\href{https://unimore.coursecatalogue.cineca.it/insegnamenti/2024/27831/2024/10001/10858?coorte=2024}{Collegamento al catalogo universitario degli insegnamenti}}

\cventry{2020 - presente}{``Advanced Sensors for Electric Vehicles''}
{Laurea Magistrale in Electric Vehicle Engineering, I anno - 1° semestre}
{\newline  Ingegneria dell'Energia Elettrica e dell'Informazione ``Guglielmo Marconi'', Università degli Studi di Bologna}{\newline \textbf{Responsabilità dell'Insegnamento Integrato di 6 CFU - Titolare di Modulo di 3 CFU (30 ore, IMIS-01/B - ex ING-INF/07)}\newline
    L'insegnamento è strutturato in due moduli di 3 CFU ciascuno. Il candidato è responsabile dell'intero insegnamento integrato, e titolare di uno dei due moduli.\newline 5 anni, dal A.A. 2020/2021}{\href{https://www.unibo.it/en/study/course-units-transferable-skills-moocs/course-unit-catalogue?codiceMateria=93730&annoAccademico=2024&codiceCorso=5699&single=True&search=True}{Collegamento al catalogo universitario degli insegnamenti} - \href{https://www.motorvehicleuniversity.com/master-degree/electric-vehicle-engineering/}{https://www.motorvehicleuniversity.com/master-degree/electric-vehicle-engineering/}}


\cventry{2020 - 2023}{``Misure Elettroniche e Laboratorio''}
{Laurea Triennale in Ingegneria Elettronica, III anno - 2° semestre}
{\newline  Dip. di Ingegneria “Enzo Ferrari'', Università degli Studi di Modena}{\newline \textbf{Responsabilità di Modulo 2 CFU (30 ore, IMIS-01/B - ex ING-INF/07)} - Titolarità del corso: prof. L. Rovati (SSD: ING-INF/07 e ING-INF/01)\newline
    L'insegnamento è strutturato in 6 CFU di lezioni frontali e 3 CFU di laboratorio. Il laboratorio prevede una parte di allineamento all'utilizzo di LabVIEW e una parte di sviluppo autonomo di un'attività progettuale per la realizzazione sistema di misura, basato sulle tecniche e i sensori descritti a lezione.\newline 3 anni, dal A.A. 2020/2021 al A.A. 2022/2023}{\href{https://unimore.coursecatalogue.cineca.it/insegnamenti/2022/24151/2016/10000/10359?coorte=2020&schemaid=19254}{Collegamento al catalogo universitario degli insegnamenti}}

\subsection{Incarichi di Insegnamento precedenti alla nomina a Professore di II fascia}

\cventry{2019 - 2020}{``Controllori e Reti di Comunicazioni Industriali''}
{Laurea Triennale in Ingegneria Meccatronica., III anno - 2° semestre}
{\newline  Dipartimento di Tecnica e Gestione dei Sistemi Industriali, Università degli Studi di Padova\newline \textbf{Titolarità dell'Insegnamento - 6 CFU}\newline
    L'insegnamento era strutturato in 4 CFU di lezioni frontali e 2 CFU di laboratorio. L'attività di laboratorio era diretta all'apprendimento dei fondamenti sulla programmazione dei PLC e dell'interfacciamento con i sensori/attuatori ad esso collegati}{\newline 1 anno, A.A. 2019/2020}{\href{https://didattica.unipd.it/off/2017/LT/IN/IN2376/000ZZ/INP7078641/N0}{Collegamento al catalogo universitario degli insegnamenti}}

\cventry{2015 - 2018}{``Reti di Comunicazioni Industriali''}
{Laurea Triennale in Ingegneria Meccanica e Meccatronica., III anno - 2° semestre}
{\newline  Dipartimento di Tecnica e Gestione dei Sistemi Industriali, Università degli Studi di Padova\newline \textbf{Titolarità dell'Insegnamento - 6 CFU}\newline
    L'insegnamento  era strutturato in 4 CFU di lezioni frontali e 2 CFU di laboratorio. L'attività di laboratorio era diretta all'apprendimento dei fondamenti sulla programmazione dei PLC e dell'interfacciamento con i sensori/attuatori ad esso collegati}{\newline 3 anni, dal A.A. 2015/2016 al A.A. 2017/2018}{\href{https://didattica.unipd.it/off/2015/LT/IN/IN0516/000ZZ/INP3058418/N0}{Collegamento al catalogo universitario degli insegnamenti}}


\cventry{2018 - 2019}{``Fondamenti di Informatica''}
    {Laurea Triennale in Ingegneria Meccatronica., I anno - 1° semestre}
    {\newline  Dipartimento di Tecnica e Gestione dei Sistemi Industriali, Università degli Studi di Padova\newline \textbf{Titolarità dell'Insegnamento - 9 CFU}\newline
        L'insegnamento  era strutturato in 6 CFU di lezioni frontali e 3 CFU di laboratorio. Come insegnamento del primo ano della triennale, la numerosità degli studenti era di circa 300 discenti. L'attività di laboratorio era diretta all'apprendimento dei fondamenti sulla programmazione in C.}{\newline 1 anno, A.A. 2018/2019}{\href{https://didattica.unipd.it/off/2018/LT/IN/IN2376/000ZZ/IN18103361/N0}{Collegamento al catalogo universitario degli insegnamenti}}

\cventry{2016 - 2018}{``Programmazione di Sistemi Embedded''}
    {Laurea Magistrale in Ingegneria Meccatronica., I anno - 1° semestre}
    {\newline  Dipartimento di Tecnica e Gestione dei Sistemi Industriali, Università degli Studi di Padova\newline \textbf{Titolarità dell'Insegnamento - 9 CFU}\newline
        L'insegnamento era rivolto agli studenti del primo anno della Laurea magistrale in Meccatronica, ed era rivolto alle nozioni teoriche a alle basi di programmazione dei sistemi embedded real-time.}{\newline 2 anni, dal A.A. 2016/2017 al 2017/2018}{\href{https://didattica.unipd.it/off/2017/LM/IN/IN0529/000ZZ/IN01122661/N0}{Collegamento al catalogo universitario degli insegnamenti}}


\cventry{A.A. 2012/2013\newline II semestre}{Lezioni per l'Insegnamento di ``Compatibilità Elettromagnetica'', 6 ore}
{Titolare: prof. M. Bertocco\newline Laurea Triennale in Ingegneria Biomedica.\newline
Dipartimento di Ingegneria dell'Informazione}
{Università degli Studi di Padova}{}{}

\cventry{A.A. 2011/2012\newline II semestre}{Lezioni per l'Insegnamento di ``Compatibilità Elettromagnetica'', 6 ore}
{Titolare: prof. M. Bertocco\newline Laurea Triennale in Ingegneria Elettronica.\newline
Dipartimento di Ingegneria dell'Informazione}
{Università degli Studi di Padova}{}{}
}

%%%%%%% INGLESE

\langen{

\cvitem{}{The current institutional teaching load is 21 ECTS (182 hours), divided into 5 courses across three universities (University of Modena and Reggio Emilia, University of Bologna, University of Verona).\medskip
    \begin{itemize}
        \item 6 ECTS (48 hours) of ``Methods, Measurement Techniques, and Sensors'' (Interuniversity Bachelor's Degree in Medical Systems Engineering for the Person, UniVR)
        \item 1 ECTS (16 hours) of ``Circuit Electronics Laboratory'' (Bachelor's Degree in Electronic Engineering, Unimore)
        \item 9 ECTS (72 hours) of ``Industrial Measurements'' (Master's Degree in Electronics Engineering, Unimore)
        \item 3 ECTS (30 hours) of ``Advanced Sensors for Electric Vehicles'' (Master's Degree in Electric Vehicle Engineering, UniBO)
        \item 2 ECTS (16 hours) of ``Hardware and Software Co-Design'' (Master's Degree in Electronics Engineering, Unimore)
    \end{itemize}
    \mbox{}\newline
    Additionally, there is the responsibility for a course, totaling 10 hours of frontal teaching, introduced in the current academic year, as part of the ``Teacher Training Programs for Secondary Schools'' conducted at the University of Modena.\medskip\newline
    Below are the details of the courses (or modules) for which there is responsibility, as well as information on courses held prior to the appointment as Associate Professor.}\medskip


\cventry{2024 - present}{``Hardware and Software Co-Design''}
{Master's Degree in Electronics Engineering, 1st Year - 1st Semester}
{\newline  Department of Engineering “Enzo Ferrari'', University of Modena\newline \textbf{Module Responsibility: 2 ECTS (16 hours, IMIS-01/B - formerly ING-INF/07)} - Course Coordinator: Prof. L. Rovati (SSD: ING-INF/07 and ING-INF/01)\newline
    The course, a core subject for the first year of the Master's in Electronics Engineering, totaling 12 ECTS, is divided into two modules of 6 ECTS each, the first in the electronics area (ING-INF/01) and the second in the electronic measurements area (ING-INF/07). This second module of 6 ECTS is further divided into three submodules of 2 ECTS each. In the 2 ECTS module under responsibility, students are introduced to the simulation of simple electronic circuits using SPICE, with particular reference to signal filtering stages useful in signal conditioning.}{\newline 1 year, starting A.Y. 2024/2025}{\href{https://unimore.coursecatalogue.cineca.it/insegnamenti/2024/28528/2024/10001/10858?coorte=2024&schemaid=21692}{Link to the university course catalog}}

\cventry{2025 - }{``Signal Generation, Acquisition, Processing, and Distribution''}{Teacher Training Programs for Secondary Schools}
{\newline  Department of Engineering “Enzo Ferrari'', University of Modena\newline \textbf{Course Responsibility: 10 hours (SSD: ING-INF/07)}\newline
    The course, consisting of 10 hours of frontal teaching, is part of the teacher training program for secondary schools, conducted in-person, and is directed at teaching classes A040 (Electrical and Electronic Science and Technology) and B015 (Laboratories of Electrical and Electronic Science and Technology).}{\newline 1 year, A.Y. 2024/2025}{\href{https://www.unimore.it/it/didattica/formazione-insegnanti/percorsi-formativi-abilitanti-60-cfu-la-scuola-secondaria}{https://www.unimore.it/it/didattica/formazione-insegnanti/percorsi-formativi-abilitanti-60-cfu-la-scuola-secondaria}}

\cventry{2023 - 2024}{``Connections and Networks, RS-232, LAN, Fundamentals of Field Buses''}{Teacher Training Programs for Secondary Schools}
{\newline  Department of Engineering “Enzo Ferrari'', University of Modena\newline \textbf{Course Responsibility: 10 hours (SSD: ING-INF/07)}\newline
    The course, consisting of 10 hours of frontal teaching, is part of the teacher training program for secondary schools, conducted in a blended mode, and is directed at teaching class B015 (Laboratories of Electrical and Electronic Science and Technology).}{\newline 1 year, A.Y. 2023/2024}{\href{https://www.unimore.it/it/didattica/formazione-insegnanti/percorsi-formativi-abilitanti-60-cfu-la-scuola-secondaria}{https://www.unimore.it/it/didattica/formazione-insegnanti/percorsi-formativi-abilitanti-60-cfu-la-scuola-secondaria}}

\cventry{2023 - present}{``Methods, Measurement Techniques, and Sensors''}
{Interuniversity Bachelor's Degree in Medical Systems Engineering for the Person, 3rd Year - 1st Semester}
{\newline  Department of Engineering for Innovation in Medicine, University of Verona}{\newline \textbf{Course Responsibility: 6 ECTS (48 hours, IMIS-01/B - formerly ING-INF/07)} -\newline
    The course is part of an interuniversity degree program among the universities of Verona, Modena and Reggio Emilia, and Trento. The administrative headquarters of the program is at the University of Verona, where the lectures are held.\newline 2 years, starting A.Y. 2023/2024}{\href{https://www.corsi.univr.it/?ent=cs&aa=2024\%2F2025&codiceCs=S25&codins=4S009924&discr=&discrCd=&id=1001&menu=studiare&tab=insegnamenti}{Link to the university course catalog}}

\cventry{2023 - present}{``Circuit Electronics Laboratory''}
{Bachelor's Degree in Electronic Engineering, 1st Year - 2nd Semester}
{\newline  Department of Engineering “Enzo Ferrari'', University of Modena}{\newline \textbf{Module Responsibility: 1 ECTS (16 hours, IMIS-01/B - formerly ING-INF/07)} - Course Coordinator: Prof. L. Vincetti (formerly ING-INF/01)\newline
    The course, strongly oriented towards laboratory activities, is structured into 6 ECTS divided into modules of 1 ECTS, corresponding to 16 hours of laboratory activities, due to the need to divide students into consecutive shifts. Each module is taught by different professors, belonging to the fields of electronics and electronic measurements. The module of 1 ECTS under responsibility focuses on the analysis and design of electronic circuits, with particular attention to the design and realization of a measurement system.\newline 2 years, starting A.Y. 2023/2024}{\href{https://unimore.coursecatalogue.cineca.it/insegnamenti/2024/27909/2021/10000/10359?coorte=2024}{Link to the university course catalog}}

\cventry{2020 - present}{``Industrial Measurements''}
{Master's Degree in Electronics Engineering, 1st Year - 2nd Semester}
{\newline  Department of Engineering “Enzo Ferrari'', University of Modena}{\newline \textbf{Course Responsibility: 9 ECTS (72 hours, IMIS-01/B - formerly ING-INF/07)}\newline
    The course is structured into 6 ECTS of lectures and 3 ECTS of laboratory work. The laboratory includes an introduction to the use of LabVIEW and an independent project activity for the realization of a measurement system, based on the techniques and sensors described in the lectures.\newline 5 years, starting A.Y. 2020/2021}{\href{https://unimore.coursecatalogue.cineca.it/insegnamenti/2024/27831/2024/10001/10858?coorte=2024}{Link to the university course catalog}}

\cventry{2020 - present}{``Advanced Sensors for Electric Vehicles''}
{Master's Degree in Electric Vehicle Engineering, 1st Year - 1st Semester}
{\newline  Department of Electrical, Electronic, and Information Engineering ``Guglielmo Marconi'', University of Bologna}{\newline \textbf{Responsibility for the Integrated Course of 6 ECTS - Module Coordinator for 3 ECTS (30 hours, IMIS-01/B - formerly ING-INF/07)}\newline
    The course is structured into two modules of 3 ECTS each. The candidate is responsible for the entire integrated course and is the coordinator of one of the two modules.\newline 5 years, starting A.Y. 2020/2021}{\href{https://www.unibo.it/en/study/course-units-transferable-skills-moocs/course-unit-catalogue?codiceMateria=93730&annoAccademico=2024&codiceCorso=5699&single=True&search=True}{Link to the university course catalog} - \href{https://www.motorvehicleuniversity.com/master-degree/electric-vehicle-engineering/}{https://www.motorvehicleuniversity.com/master-degree/electric-vehicle-engineering/}}


\cventry{2020 - 2023}{``Electronic Measurements and Laboratory''}
{Bachelor's Degree in Electronic Engineering, 3rd Year - 2nd Semester}
{\newline  Department of Engineering “Enzo Ferrari'', University of Modena}{\newline \textbf{Module Responsibility: 2 ECTS (30 hours, IMIS-01/B - formerly ING-INF/07)} - Course Coordinator: Prof. L. Rovati (SSD: ING-INF/07 and ING-INF/01)\newline
    The course is structured into 6 ECTS of lectures and 3 ECTS of laboratory work. The laboratory includes an introduction to the use of LabVIEW and an independent project activity for the realization of a measurement system, based on the techniques and sensors described in the lectures.\newline 3 years, from A.Y. 2020/2021 to A.Y. 2022/2023}{\href{https://unimore.coursecatalogue.cineca.it/insegnamenti/2022/24151/2016/10000/10359?coorte=2020&schemaid=19254}{Link to the university course catalog}}

\subsection{Teaching Assignments Prior to Appointment as Associate Professor}

\cventry{2019 - 2020}{``Controllers and Industrial Communication Networks''}
{Bachelor's Degree in Mechatronic Engineering, 3rd Year - 2nd Semester}
{\newline  Department of Management and Engineering, University of Padova\newline \textbf{Course Responsibility: 6 ECTS}\newline
    The course was structured into 4 ECTS of lectures and 2 ECTS of laboratory work. The laboratory activities focused on the fundamentals of PLC programming and interfacing with connected sensors/actuators.}{\newline 1 year, A.Y. 2019/2020}{\href{https://didattica.unipd.it/off/2017/LT/IN/IN2376/000ZZ/INP7078641/N0}{Link to the university course catalog}}

\cventry{2015 - 2018}{``Industrial Communication Networks''}
{Bachelor's Degree in Mechanical and Mechatronic Engineering, 3rd Year - 2nd Semester}
{\newline  Department of Management and Engineering, University of Padova\newline \textbf{Course Responsibility: 6 ECTS}\newline
    The course was structured into 4 ECTS of lectures and 2 ECTS of laboratory work. The laboratory activities focused on the fundamentals of PLC programming and interfacing with connected sensors/actuators.}{\newline 3 years, from A.Y. 2015/2016 to A.Y. 2017/2018}{\href{https://didattica.unipd.it/off/2015/LT/IN/IN0516/000ZZ/INP3058418/N0}{Link to the university course catalog}}


\cventry{2018 - 2019}{``Foundations of Computer Science''}
    {Bachelor's Degree in Mechatronic Engineering, 1st Year - 1st Semester}
    {\newline  Department of Management and Engineering, University of Padova\newline \textbf{Course Responsibility: 9 ECTS}\newline
        The course was structured into 6 ECTS of lectures and 3 ECTS of laboratory work. As a first-year course in the Bachelor's program, the student cohort was approximately 300. The laboratory activities focused on the fundamentals of programming in C.}{\newline 1 year, A.Y. 2018/2019}{\href{https://didattica.unipd.it/off/2018/LT/IN/IN2376/000ZZ/IN18103361/N0}{Link to the university course catalog}}

\cventry{2016 - 2018}{``Embedded Systems Programming''}
    {Master's Degree in Mechatronic Engineering, 1st Year - 1st Semester}
    {\newline  Department of Management and Engineering, University of Padova\newline \textbf{Course Responsibility: 9 ECTS}\newline
        The course was aimed at first-year Master's students in Mechatronics and focused on theoretical concepts and the basics of real-time embedded systems programming.}{\newline 2 years, from A.Y. 2016/2017 to 2017/2018}{\href{https://didattica.unipd.it/off/2017/LM/IN/IN0529/000ZZ/IN01122661/N0}{Link to the university course catalog}}



% \cventry{A.A. 2017/2018}{``Embedded Systems Programming'', 9 CFU}
% {Responsibility of the course.\newline Master Degree in Mechatronic Engineering.
%     Department of Management and Engineering  (DTG)}
% {University of Padova}{}
% {S. S. Nicola, 3, \texttt{I-36100} --- Vicenza, VI}

% \cventry{A.A. 2016/2017}{``Embedded Systems Programming'', 9 CFU}
% {Responsibility of the course.\newline Master Degree in Mechatronic Engineering.
%     Department of Management and Engineering  (DTG)}
% {University of Padova}{}
% {S. S. Nicola, 3, \texttt{I-36100} --- Vicenza, VI}

% \cventry{A.A. 2016/2017}{``Industrial Communication Networks'', 6 CFU}
% {Responsibility of the course.\newline Bachelor Degree in Mechanics and Mechatronic Engineering.
%     Department of Management and Engineering  (DTG)}
% {University of Padova}{}
% {S. S. Nicola, 3, \texttt{I-36100} --- Vicenza, VI}

% \cventry{A.A. 2015/2016}{``Industrial Communication Networks'', 6 CFU}
% {Responsibility of the course.\newline Bachelor Degree in Mechanics and Mechatronic Engineering.
%     Department of Management and Engineering  (DTG)}
% {University of Padova}{}
% {S. S. Nicola, 3, \texttt{I-36100} --- Vicenza, VI}

% \cventry{A.A. 2017/2018}{``Industrial Communication Networks'', 6 CFU}
% {Responsibility of the course.\newline Bachelor Degree in Mechanics and Mechatronic Engineering.
%     Department of Management and Engineering  (DTG)}
% {University of Padova}{New Assignment, relevant to the call D310000-1016393-2017 of the University of Padova, June 15th 2017.}
% {S. S. Nicola, 3, \texttt{I-36100} --- Vicenza, VI}

}

\subsectioniten{Relatore di Tesi di Laurea}{Thesis Supervisor}

\cvitem{2025}{``Testing and validation of a new generation map-based Adaptive Cruise Control'', A. Bellei, Laurea Magistrale in Electric Vehicle Engineering. Università di Bologna. Durata tesi: 9 mesi. Data discussione: 24 marzo 2025.}

\cvitem{2024}{``Integrazione di pannello fotovoltaico e batteria per un sistema di monitoraggio ininterrotto delle emissioni di vapore nel sottosuolo ​'', D. Parrella, Laurea Triennale in Ingegneria Elettronica. Università di Modena e Reggio Emilia. Durata tesi: 4 mesi. Data discussione: 9 luglio 2024.}

\cvitem{2024}{``Range and Low-Voltage Accessory Current Testing of the E260 Electric Vehicle in Various Environmental Conditions'', Jingrui Luo, Laurea Magistrale in Electric Vehicle Engineering. Università di Bologna. Durata tesi: 9 mesi. Data discussione: 22 luglio 2024.}

\cvitem{2020}{``Protocollo OPC UA per l'interoperabilità nei sistemi industriali'', G. Cannavale, Laurea Triennale in Ingegneria Meccatronica. Durata tesi: 4 mesi.}

\cvitem{2018}{``OPC UA e TSN. Analisi di una soluzione innovativa
per le comunicazioni industriali'', S. Vanti, Laurea Triennale in Ingegneria Meccatronica. Università di Padova. Durata tesi: 4 mesi. Data discussione: 16 luglio 2018.}

\cvitem{2017}{``Analisi delle prestazioni di un link wireless IEEE 802.11 per l’interconnessione tra due segmenti real-time ethernet industriali'', E. Petri, Laurea Triennale in Ingegneria Meccanica e Meccatronica. Università di Padova. Durata tesi: 4 mesi. Data discussione: 28 novembre 2017.}

\cvitem{2017}{``Programmazione di Dispositivi Industriali per la misura delle prestazioni di reti real-time Profibus DP'', R. Salviati, Laurea Triennale in Ingegneria Meccanica e Meccatronica. Università di Padova. Durata tesi: 4 mesi. Data discussione: 28 novembre 2017.}

\cvitem{2017}{``Programmazione di Dispositivi Industriali per il Collegamento a Reti Real-Time Profibus DP'', M. Simioni, Laurea Triennale in Ingegneria Meccanica e Meccatronica. Università di Padova. Durata tesi: 4 mesi. Data discussione: 18 settembre 2017.}

\subsection{Correlatore di Tesi di Laurea}

\cvitem{2016}{``Analysis of Ethernet Powerlink network and development of a wireless extension based on the IEEE 802.11n WLAN'', A. Tagliapietra, Laurea Magistrale in Ingegneria delle Telecomunicazioni. Università di Padova. 
Durata tesi: 6 mesi. Data discussione: 14 dicembre 2016.}

\cvitem{2016}{``Misure di Consumo energetico per schede Ethernet conformi allo standard EEE'', R. Mollaiyan, Laurea Triennale in Ingegneria Elettronica. DUniversità di Padova. urata tesi: 4 mesi. Data discussione: 23 Luglio 2016.}

\cvitem{2015}{``Stima di Canale per Autenticazione mediante 
Software Defined Radio'', A. Cavicchio, Laurea Triennale in Ingegneria Informatica. Università di Padova. 
Durata tesi: 9 mesi. Data discussione: 18 Marzo 2015.}

\cvitem{2015}{``Powerline Communication per sistemi
Domotici'', M. Dominese, Laurea Triennale in Ingegneria dell'Automazione. Università di Padova. 
Durata tesi: 3 mesi. Data discussione: 18 Marzo 2015.}

\cvitem{2014}{``Analysis of an IEEE 802.11-based protocol for real-time applications in agriculture'', M. Luvisotto, Laurea Magistrale in Ingegneria dell'Automazione. Università di Padova. 
Durata tesi: 6 mesi. Data discussione: 7 Luglio 2014.}

\cvitem{2014}{``Implementazione dello standard IEEE 802.11g su 
sistemi Software Defined Radio (SDR)'', 
F. Casamento, Laurea Triennale in Ingegneria dell'Informazione. Università di Padova. 
Durata tesi: 3 mesi. Data discussione: 20 febbraio 2014}


\cvitem{2013}{``Traffic Models for Data Networks and Energy 
Efficient Control'', M. Michielan, Laurea Magistrale in Ingegneria dell'Au-tomazione. Università di Padova. 
Durata tesi: 7 mesi. Data discussione: 
9 Dicembre 2013}

\cvitem{2013}{``Reti ethernet a efficienza energetica per applicazioni industriali. Parte A: lo standard 802.3az'', 
M. Conti, Laurea Triennale in Ingegneria dell'Informazione. Università di Padova. 
Durata tesi: 4 mesi. Data discussione: 30 settembre 2013}

\cvitem{2013}{``Reti ethernet a efficienza energetica per applicazioni industriali. Parte B: le misure in laboratorio'', 
M. Perazzolo, Laurea Triennale in Ingegneria dell'Informazione. Università di Padova. 
Durata tesi: 4 mesi. Data discussione: 30 settembre 2013}

\cvitem{2013}{``Ottimizzazione del consumo energetico per reti
Ethernet Real-Time'', A. Salata, Laurea
Magistrale in Ingegneria dell'Au-tomazione. Università di Padova. 
Durata tesi: 6 mesi. Data discussione: 11 Marzo 2013}

\cvitem{2012}{``Studio e progetto di metodi di misura robusti
per sincrofasori in bassa tensione'', M. Penzo, Laurea
Magistrale in Ingegneria Elettronica. Università di Padova. 
Durata tesi: 7 mesi. Data discussione: Ottobre 2012}

\cvitem{2013}{`Impiego di un dispositivo basato su FPGA per misure
di power quality in Real-Time'', I. Ergest, Laurea
Triennale in Ingegneria Elettronica. Università di Padova. 
Durata tesi: 5 mesi. Data discussione: 25 Febbraio 2013}

\cvitem{2011}{``Simulazione cross-layer di tecniche di ritrasmissione
in reti wireless interferenti 802.11 e 802.15.4'', R. Bersan, Laurea
Triennale in Ingegneria Elettronica. Università di Padova. 
Durata tesi: 4 mesi. Data discussione: 2011}